\section{System Model and Implemented Methods}
\label{sec:method}
\subsection{Two Distinct Evidence Paths}
The benchmark has two paths (Fig.~\ref{fig:architecture}). The first reclassifies stored service-profile snapshots and applies a context-name mapping. The second produces seeded abstract timelines from synthetic class priors. They share outcome terminology, but neither path supplies the other's observations. In particular, the timeline does not derive its Bernoulli probabilities from the profile corpus, and its class names differ from the profile names. This distinction prevents an apparent combined validation from two artifacts that were never coupled.

\begin{figure}[t]
\centering
\fbox{\begin{minipage}{0.91\columnwidth}\small
\textbf{Path A: semantic corpus}\par
Frozen profiles + workload coverage $\longrightarrow$\par
State classification $\longrightarrow$ 21 labels\par
Authored context map $\longrightarrow$ 20 overlay rows\medskip\par
\textbf{Path B: stochastic generator}\par
Seed + four synthetic priors $\longrightarrow$\par
48 abstract steps per class $\longrightarrow$ $C_s,\tau_s$\par
30 seed means $\longrightarrow$ Student-$t$ intervals\medskip\par
\textbf{Audit layer}\par
Frozen hashes, source parity, metric dependencies,\par
post hoc analytic expectations and limitations
\end{minipage}}
\caption{Implemented evidence flow. No live sensors, channel calibration, workload execution, or policy optimizer connects the two paths.}\label{fig:architecture}
\end{figure}

\subsection{Profile-Based Minimum Useful Service}
Let $o$ be the stored orchestrator state and $F_{c,w}$ a Boolean declaration that profile $c$ supports workload $w$ offline. The implemented classifier is
\begin{equation}
L(o,F)=\begin{cases}
3,&o=\mathrm{connected},\\
2,&o\in\{\mathrm{degraded},\mathrm{transitioning}\},\\
1,&o=\mathrm{offline},\ F=1,\\
0,&\text{otherwise}.
\end{cases}\label{eq:classifier}
\end{equation}
The labels are target, degraded, minimum useful, and failed, respectively. The semantic acceptance predicate is $L\geq1$. These are ordered categories, not equally spaced utility scores. Averaging their numeric encodings would require a separate cardinal-utility assumption and is not done here.

For a hand check, an offline local-creation snapshot with declared workload support receives $L=1$, whereas an offline wearable coding request without support receives $L=0$. Changing only the stored state to connected yields $L=3$ for either request. This last case exposes a classifier limitation: capability and application success are not evaluated on the connected branch. These three evaluations illustrate the rule; they are not added observations.

The four profiles are desk, mobile/docked, local creation, and wearable. Each contains five common scenario snapshots; wearable includes a sixth unsupported coding request. Connected state receives the target label without checking measured latency, throughput, task correctness, or resource availability. Thus, \eqref{eq:classifier} operationalizes the repository's current \emph{declared} service threshold, not a user-validated threshold. It can expose an unsupported offline workload, but cannot detect a connected application failure if the stored state remains connected.

Five Gary context names map to common snapshots: neighborhood outage maps to jam/offline fallback; storm congestion, school bandwidth collapse, and library hotspot overload map to degraded Wi-Fi; payment-system unavailability maps to forced offline. The mapping is fixed in the versioned configuration. It is a lookup assumption, not a model of those events' physical causes or business consequences.

\subsection{Seeded Continuity Generator}
For each seed $s$, the simulator instantiates \texttt{random.Random(s)} and iterates over wearable, smartphone, laptop, and edge node, each for $T=48$ abstract steps. Let $a_c$, $q_c$, and $h_c$ denote class useful-service base probability, conditional task base probability, and stress sensitivity. Table~\ref{tab:priors} gives their exact values. At each step the code draws stress $S\sim U(0,1)$ and demand $D\sim U(1,50)$, then samples
\begin{align}
M_{sct}&\sim\mathrm{Bernoulli}\big[a_c(1-0.35h_cS)\big],\label{eq:useful}\\
Y_{sct}&=M_{sct}B_{sct},\label{eq:task}\\
B_{sct}&\sim\mathrm{Bernoulli}\big[q_c(1-0.2S)\big].\notag
\end{align}
The task draw is performed only if $M_{sct}=1$, exactly as in the implementation. Stress and Bernoulli draws use a pseudorandom stream; the displayed probability model describes their intended sampling semantics. All configured probabilities remain within $[0,1]$.

\begin{table}[t]\centering\caption{Synthetic timeline priors. Values are authored assumptions, not calibrated device measurements.}\label{tab:priors}\small
\begin{tabular}{lrrr}\toprule
Class & $a_c$ & $q_c$ & $h_c$\\\midrule
Wearable & 0.45 & 0.55 & 1.35\\
Smartphone & 0.72 & 0.78 & 1.00\\
Laptop & 0.80 & 0.85 & 0.85\\
Edge node & 0.88 & 0.90 & 0.70\\\bottomrule
\end{tabular}\end{table}

Useful occupancy and task completion for a class and seed are
\begin{equation}
\tau_{sc}=\frac1T\sum_{t=1}^{T}M_{sct},\qquad
C_{sc}=\frac1T\sum_{t=1}^{T}Y_{sct}.\label{eq:outcomes}
\end{equation}
The primary seed means weight all four classes equally: $\tau_s=\frac14\sum_c\tau_{sc}$ and $C_s=\frac14\sum_c C_{sc}$. Each step represents one synthetic task opportunity; it has no assigned number of seconds. Consequently, useful time is a fraction of equally weighted abstract steps, and completion is a synthetic success fraction rather than a measured jobs-per-second rate. There are 5,760 class-step outcomes across the 30 seeds, but the statistical sample size is 30 seed means.

Demand does not enter \eqref{eq:useful} or \eqref{eq:task}. It only contributes to the energy proxy. There is no queue, battery state, memory constraint, congestion feedback, traffic scheduler, radio channel, or admission-control decision in this generator. Resource constraint is represented indirectly by class-specific probabilities. A class ordering in the results therefore does not demonstrate a device performance ordering under equal physical workloads.

\subsection{Reporting Algorithm and Numerical Conventions}
The implementation loads the frozen bundle, rederives all stored labels, constructs the context overlay, generates a timeline for each recorded seed, and reports primary and secondary summaries. These steps are deterministic for the source, configuration, and Python random generator used. Reproduction must preserve class order and conditional draw behavior: consuming an extra random number, even for an unused task draw, changes subsequent samples.

Class fractions and seed aggregates are rounded to six decimals before the stored report computes across-seed summaries. The manuscript reports four decimals. The analytic model applies to the unrounded fractions. The maximum rounding discrepancy between an unrounded overall seed fraction and its stored value is at most $10^{-6}$, combining class and aggregate rounding. The statistical audit independently recalculates the stored means and intervals; it does not silently replace historical values with unrounded recalculations.
